\documentclass[11pt]{article}
\usepackage[margin=.9in]{geometry}
\usepackage{amsmath,amssymb,amsthm}
\usepackage[T1]{fontenc}
\usepackage{lmodern,microtype,hyperref}
\newtheorem{theorem}{Theorem}
\newtheorem{lemma}[theorem]{Lemma}
\newtheorem{proposition}[theorem]{Proposition}
\newtheorem{corollary}[theorem]{Corollary}
\theoremstyle{remark}\newtheorem{remark}[theorem]{Remark}
\newcommand{\norm}[1]{\left\lVert#1\right\rVert}
\DeclareMathOperator{\diag}{diag}\DeclareMathOperator{\ran}{ran}
\newcommand{\R}{\mathbb R}\newcommand{\C}{\mathbb C}
\begin{document}
% No preamble. Notebook 15 minimizer module. Unique labels: nm-.
% Requires amsmath, amssymb, amsthm and a one-argument \norm macro.

\section{A uniform local constrained Wilson minimum on the flat gauge slice}
\label{nm-section}

Fix an odd $L\ge3$, $M=L^j$, $\eta=M^{-1}$, final coarse scale one,
and $0<\alpha<1$. The primary statement is on finite compatible
periodic regulators. Constants may depend on $L$, dimension, the
color norm, and $\alpha$, but not on $j$ or periodic volume. All
transposes use the physical fine pairing $\eta^d\sum\tau(AB)$ and
the corresponding coarse pairing; at complex fields they are
bilinear transposes. Write the corrected nonlinear average in its
fixed-reference potential coordinate as
\begin{equation}\label{nm-average}
 \mathcal F(A)=QA+\mathcal R(A),\qquad Q=Q_M^{\rm st},\qquad
 \mathcal R(0)=D\mathcal R(0)=0.
\end{equation}
The symbol $R$ without calligraphic font continues to denote the
fixed flat scalar gauge projection. In particular, $R$ and
$\mathcal R$ are different objects.

\subsection{Two uniform estimates for the nonlinear constraint}
\label{nm-average-estimates}

The previously established nonlinear averaging ball gives a
one-step dimensionless corrected map $F_L$ and physical map
\begin{equation}\label{nm-one-step}
 \mathcal F_h(A)=(Lh)^{-1}F_L(hA)=Q_hA+\mathcal R_h(A).
\end{equation}
Let $m_L$ bound the number of input bonds read by one output, and
$N_L$ the number of outputs reading one input. On a fixed local
complex neighborhood choose $K_0$ so that each matrix entry of
$D(F_L-LQ_h)(X)$ has norm at most $K_0\norm{X}_\infty$.
The corresponding remainder entries in physical coordinates are
bounded by $K_0h\norm{A}_\infty/L$. With
\begin{equation}\label{nm-local-constants}
 D_\infty=\frac{m_LK_0}{L},\qquad
 D_2=K_0\sqrt{m_LN_L}\,L^{d/2-1},\qquad
 D_*=N_LK_0L^{d-1},\qquad S=\frac1{L-1},
\end{equation}
the physical transpose therefore satisfies
\begin{equation}\label{nm-adjoint-one-step}
 \norm{D\mathcal R_h(A)^T}_{\infty\to\infty}
 \le D_*h\norm{A}_\infty.
\end{equation}
Indeed the transpose converts an ordinary column sum to a physical
row sum with factor $(Lh)^d/h^d=L^d$. This factor is already
included in $D_*$.

Both $Q_h$ and its physical transpose are bounded-field
contractions. If $a=\norm{A}_\infty$ belongs to the common nonlinear
ball, the intermediate fields obey $\norm{A_t}_\infty\le2a$.
The spacings $h_t=L^{t-j}$ satisfy $\sum_{t<j}h_t\le S$.
Expanding the difference of the transposed nonlinear derivative
product and its flat product thus proves
\begin{equation}\label{nm-adjoint-composite}
 \norm{D\mathcal R(A)^T}_{\infty\to\infty}
 \le e^{2D_*aS}-1\le c_*a
\end{equation}
on a smaller fixed ball. The analogous derivative bounds in the
bounded and physical Hilbert norms follow with $D_\infty$ and
$D_2$. These are absolute physical row or column estimates, not
inferences from the fine dimension. The local analytic expansions
also make $A\mapsto D\mathcal R(A)^T$ holomorphic in this operator
norm. Cauchy's estimate on a slightly larger common ball gives
\begin{equation}\label{nm-adjoint-lipschitz}
 \norm{D\mathcal R(A)^T-D\mathcal R(A')^T}_{\infty\to\infty}
 \le c_{*,1}\norm{A-A'}_\infty.
\end{equation}
No bounded-field estimate for the gauge projection $\Pi$ is used.

We also need a bound with two Hilbert directions, rather than just
two bounded directions. Let $K_2$ bound the local bilinear Hessian
$D^2F_L$ in the maximum norms of its inputs; it bounds the full
second derivative, including its factor two. The scaling
$D^2\mathcal F_h=(h/L)D^2F_L(hA)$, local Cauchy--Schwarz, and
incidence counting give
\begin{equation}\label{nm-hessian-one-step}
 \left|\left\langle w,D^2\mathcal F_h(A)[B,C]\right\rangle_{Lh}\right|
 \le k_2h\norm{w}_\infty\norm{B}_{2,h}\norm{C}_{2,h},
 \qquad k_2=K_2N_LL^{d-1}.
\end{equation}
For clarity, replace each local input maximum by its local
Euclidean norm. Summing the product over output bonds and applying
Cauchy--Schwarz costs at most $N_L$; restoring the fine measure
then converts $(Lh)^d(h/L)$ to the factor in
\eqref{nm-hessian-one-step}.

Let $P_t$ be the derivative of the nonlinear prefix through stage
$t$, and $T_t$ the derivative of the suffix after stage $t+1$,
both evaluated along the actual trajectory. The exact second-order
chain rule is
\begin{equation}\label{nm-hessian-chain}
 D^2\mathcal F(A)[B,C]
 =\sum_{t<j}T_tD^2\mathcal F_{h_t}(A_t)[P_tB,P_tC].
\end{equation}
The two prefix Hilbert norms are each at most $e^{2D_2aS}$ and the
suffix transpose bounded-field norm at most $e^{2D_*aS}$.
Pairing \eqref{nm-hessian-chain} with $w$ and using
\eqref{nm-hessian-one-step} therefore proves the uniform form bound
\begin{equation}\label{nm-hessian-composite}
 \boxed{\quad
 \left|\left\langle w,D^2\mathcal F(A)[B,C]\right\rangle_1\right|
 \le k_2S e^{(4D_2+2D_*)aS}
       \norm{w}_\infty\norm{B}_{2,\eta}\norm{C}_{2,\eta}.
 \quad}
\end{equation}

\subsection{The exact reduced equation}
\label{nm-reduction}

Let $K=d^Td$, let $q=q_M$ be the Maxwell coarse precision, and let
$H$ be the flat minimizing extension in the matched source gauge.
The flat identities are
\begin{equation}\label{nm-flat-identities}
 QH=I,\qquad KH=Q^Tq,\qquad RD^TH=0.
\end{equation}
Set
\begin{equation}\label{nm-slice-green}
 \mathcal X=\ker Q\cap\ker(RD^T),\qquad
 \mathcal C=\Pi\Gamma\Pi^T=\Gamma-DZD^T,
 \qquad \Gamma=G_0-H(q+a_{\rm pen})^{-1}H^T.
\end{equation}
The established constrained-inverse identities give
\begin{equation}\label{nm-constrained-identities}
 Q\mathcal C=RD^T\mathcal C=\mathcal C Q^T=0,\qquad
 \mathcal C K\xi=\xi\quad(\xi\in\mathcal X),\qquad
 \mathcal C KH=0.
\end{equation}
On a real finite regulator, $\mathcal C$ is the positive
variational inverse of $K$ on $\mathcal X$. Also
$H^TK\xi=qQ\xi=0$. The semidefinite $q$ in these identities is
not the positive auxiliary precision $q+a_{\rm pen}$.

For coarse $B$ the exact local constraint chart is
\begin{equation}\label{nm-chart}
 A=\xi+H[B-\mathcal R(A)],\qquad \xi\in\mathcal X.
\end{equation}
It gives $\mathcal F(A)=B$ and $RD^TA=0$. Conversely, every small
field satisfying these two conditions belongs to this chart,
by taking $\xi=A-H[B-\mathcal R(A)]$. Its derivative in $\xi$ is
$T_A=(I+HD\mathcal R(A))^{-1}$ restricted to $\mathcal X$.

Write the ordinary-potential Wilson Euler derivative as
\begin{equation}\label{nm-force}
 E(A)=KA+N(A),\qquad N(A)=S_{\rm jet}^Tf_\eta(S_{\rm jet}A),
\end{equation}
where $S_{\rm jet}$ collects the finitely many local values and
ordinary first differences used by the plaquette jet. Stationarity
of the action in the chart is equivalent to
$\mathcal C T_A^TE(A)=0$. The ordered inverse identity
\[
 T_A^T=I-D\mathcal R(A)^T
       (I+H^TD\mathcal R(A)^T)^{-1}H^T
\]
and \eqref{nm-flat-identities}--\eqref{nm-chart} transform this
condition exactly into
\begin{equation}\label{nm-fixed-point}
 \boxed{\quad
 \xi=-\mathcal C N(A)
 +\mathcal C D\mathcal R(A)^T
       (I+H^TD\mathcal R(A)^T)^{-1}
       \{q[B-\mathcal R(A)]+H^TN(A)\}.
 \quad}
\end{equation}
In particular, the nonlinear constraint contributes the second
term in \eqref{nm-fixed-point}; it cannot be omitted.

\subsection{Uniform analytic contraction}
\label{nm-contraction}

Use the complete-field physical norm
\begin{equation}\label{nm-norm}
 \norm{A}_{\mathcal N_\alpha}
 =\norm{A}_\infty+\norm{\nabla A}_\infty+[\nabla A]_\alpha.
\end{equation}
It controls the complete $C^\alpha$ norms of the jet, uniformly in
$\eta$. Telescoping controls differences of $A$ at distances below
one; its bounded norm handles larger distances. The flat extension
and the constrained Green theorem proved above supply
\begin{align}
 H &: \ell^\infty_c\longrightarrow\mathcal N_\alpha,&
 \mathcal C &: \ell^\infty_f\longrightarrow\mathcal N_\alpha,
                                                        \label{nm-analytic-input-a}\\
 \mathcal C S_{\rm jet}^T &: C^\alpha_f\longrightarrow\mathcal N_\alpha,
 &q &: \ell^\infty_c\longrightarrow\ell^\infty_c.
                                                        \label{nm-analytic-input-b}
\end{align}
The equal exponent in \eqref{nm-analytic-input-b} is essential.
It uses the separate constrained biharmonic cancellation, including
the fourth and fifth heat moments and the separation-dependent
time split, not merely an estimate for $G_0$ or a theorem with
loss of H\"older exponent.

The transpose estimates for $H$ and its first differences give
\begin{equation}\label{nm-transpose-input}
 H^T:\ell^\infty_f\longrightarrow\ell^\infty_c,
 \qquad H^TS_{\rm jet}^T=(S_{\rm jet}H)^T:
       C^\alpha_f\longrightarrow\ell^\infty_c.
\end{equation}
For the second map the bounded input norm already suffices:
the physical volume of a fine cell is fixed and the physical-cell
Hilbert transpose rows of $S_{\rm jet}H$ are summable. Thus
\eqref{nm-transpose-input} never differentiates a rough force
coefficient.

The Wilson jet calculation supplies a holomorphic function
$f_\eta$ with $f_\eta(0)=Df_\eta(0)=0$ and the uniform estimates
\begin{align}
 \norm{f_\eta(S_{\rm jet}A)}_{C^\alpha}
 &\le c_f\norm{A}_{\mathcal N_\alpha}^2,                 \label{nm-force-bound}\\
 \norm{f_\eta(S_{\rm jet}A)-f_\eta(S_{\rm jet}A')}_{C^\alpha}
 &\le c_f(\norm{A}_{\mathcal N_\alpha}+
           \norm{A'}_{\mathcal N_\alpha})
                 \norm{A-A'}_{\mathcal N_\alpha}.        \label{nm-force-lipschitz}
\end{align}
These estimates concern the coefficient before the outside
transpose in \eqref{nm-force}.

Shrink the common ball so that
\begin{equation}\label{nm-neumann-smallness}
 \norm{H}_{\infty\to\mathcal N_\alpha}
       \norm{D\mathcal R(A)}_{\infty\to\infty}<\frac12,
 \qquad
 \norm{H^T}_{\infty\to\infty}
       \norm{D\mathcal R(A)^T}_{\infty\to\infty}<\frac12.
\end{equation}
The first inequality makes \eqref{nm-chart} an analytic
$\mathcal N_\alpha$ chart. Put $b=\norm{B}_\infty$ and
$h_N=\norm{H}_{\infty\to\mathcal N_\alpha}$. Its contraction
estimate gives
\begin{equation}\label{nm-chart-estimates}
 \norm{A(\xi,B)}_{\mathcal N_\alpha}
 \le2(\norm{\xi}_{\mathcal N_\alpha}+h_Nb),\qquad
 \norm{A(\xi,B)-A(\xi',B)}_{\mathcal N_\alpha}
 \le2\norm{\xi-\xi'}_{\mathcal N_\alpha}.
\end{equation}
Define the coarse multiplier
\begin{equation}\label{nm-multiplier}
 m(A,B)=(I+H^TD\mathcal R(A)^T)^{-1}
             \{q[B-\mathcal R(A)]+H^TN(A)\}.
\end{equation}
For $\norm{A}_{\mathcal N_\alpha},\norm{A'}_{\mathcal N_\alpha}\le r$
in this chart, the preceding bounds imply
\begin{equation}\label{nm-multiplier-bounds}
 \norm{m(A,B)}_\infty\le C(b+r^2),\qquad
 \norm{m(A,B)-m(A',B)}_\infty
 \le C(b+r)\norm{A-A'}_{\mathcal N_\alpha}.
\end{equation}
For the first bound use inverse norm at most two,
$\norm{\mathcal R(A)}_\infty\le Cr^2$, and
$H^TN=(S_{\rm jet}H)^Tf_\eta$. For the second use the inverse
difference identity and \eqref{nm-adjoint-lipschitz}; the change
of the braces is $O(r\norm{A-A'})$, and the change of the inverse
multiplies a quantity $O(b+r^2)$.

Let $\Phi_B(\xi)$ denote the right side of
\eqref{nm-fixed-point}, with $A=A(\xi,B)$. Then
\begin{align}
 \norm{\Phi_B(\xi)}_{\mathcal N_\alpha}
 &\le C(r^2+rb+r^3),                                   \label{nm-contraction-size}\\
 \norm{\Phi_B(\xi)-\Phi_B(\xi')}_{\mathcal N_\alpha}
 &\le C(b+r)\norm{\xi-\xi'}_{\mathcal N_\alpha}.          \label{nm-contraction-lipschitz}
\end{align}
To check the constraint contribution, expand the difference of
$D\mathcal R(A)^Tm(A,B)$ into its two factors and apply
\eqref{nm-adjoint-composite}, \eqref{nm-adjoint-lipschitz}, and
\eqref{nm-multiplier-bounds}; then use the bounded-input map
$\mathcal C$ in \eqref{nm-analytic-input-a}. The force contribution
uses \eqref{nm-analytic-input-b} and
\eqref{nm-force-lipschitz}. Both contributions lie in $\mathcal X$.

Choose a fixed small $\mathcal X$-ball and then $b_*>0$ so that its
chart image stays in the preceding balls, \eqref{nm-contraction-size}
maps it into itself, and \eqref{nm-contraction-lipschitz} is below
one half. This constructs a unique small solution $\xi(B)$ for
every complex $\norm{B}_\infty<b_*$. Uniform analytic contraction
makes $\xi(B)$ and $A(B)$ holomorphic. The sharper bounds are
\begin{equation}\label{nm-quadratic-remainder}
 \norm{\xi(B)}_{\mathcal N_\alpha}\le Cb^2,
 \qquad \norm{A(B)-HB}_{\mathcal N_\alpha}\le Cb^2,
 \qquad \norm{A(B)}_{\mathcal N_\alpha}\le Cb.
\end{equation}
Indeed on $\norm{\xi}\le Kb^2$, after reducing $b_*$,
\eqref{nm-chart-estimates} gives $r\le4h_Nb$. The right side of
\eqref{nm-contraction-size} is then $C'b^2$. Take $K>2C'$, shrink
$b_*$ again so that $Kb^2\le h_Nb$ and the Lipschitz bound is below
one half, and use uniqueness in the original ball. The chart
remainder $-H\mathcal R(A)$ proves the second estimate.
If a complex coarse ball has distance $\delta$ from the boundary
of the holomorphy ball and $\norm{A(B)}_{\mathcal N_\alpha}\le C_0$
there, multivariable Cauchy gives, for $n\ge1$,
\begin{equation}\label{nm-cauchy}
 \norm{D^nA(B)[B_1,\ldots,B_n]}_{\mathcal N_\alpha}
 \le C_0 n!\left(\frac e\delta\right)^n
                    \prod_{i=1}^n\norm{B_i}_\infty.
\end{equation}
Use polyradii approaching $\delta/n$ and $n^n\le e^n n!$.

\subsection{The real solution is a strict local minimum}
\label{nm-minimum}

For real fields the ordinary-coordinate Wilson Hessian has the
form
\begin{equation}\label{nm-action-hessian}
 D^2\mathcal S(A)=K+S_{\rm jet}^TW(A)S_{\rm jet},
 \qquad \norm{W(A)}_\infty\le C\norm{A}_{\mathcal N_\alpha}.
\end{equation}
This follows from the exact jet density whose quadratic Taylor
part is the Maxwell density. It does not identify two different
coordinate Hessians at a nonzero background.

The flat full inverse's first-derivative energy bound gives
\begin{equation}\label{nm-coercivity}
 \norm{z}_{2,\eta}^2+\norm{\nabla z}_{2,\eta}^2
 \le C_E\langle z,Kz\rangle_\eta\qquad(z\in\mathcal X).
\end{equation}
In fact both the gauge and averaging penalties in the positive
full precision vanish on $\mathcal X$. Apply the uniform bound on
$(I,\nabla)G_0(I,\nabla)^T$ to its square-root factorization.

Let $z\in\mathcal X$ and $v=T_Az$. Differentiating the chart gives
$v=z-HD\mathcal R(A)v$. The cross terms in its Maxwell energy
vanish because $H^TKz=0$, while $q$ is positive semidefinite.
Consequently
\begin{equation}\label{nm-tangent-energy}
 \langle v,Kv\rangle
 =\langle z,Kz\rangle+
   \langle D\mathcal R(A)v,qD\mathcal R(A)v\rangle
 \ge\langle z,Kz\rangle.
\end{equation}
The uniform physical Hilbert bound $\norm{D\mathcal R(A)}_{2\to2}
\le Cr$, together with the Hilbert bounds for $H$ and $\nabla H$,
gives, after shrinking $r$,
\begin{equation}\label{nm-tangent-hilbert}
 \norm{v}_2\le2\norm{z}_2,\qquad
 \norm{v}_2^2+\norm{\nabla v}_2^2
 \le C\langle z,Kz\rangle.
\end{equation}
No fine Hilbert-to-pointwise bound is used here.

Differentiate \eqref{nm-chart} once more. With
$H_A=(I+HD\mathcal R(A))^{-1}H$ one obtains
$A''=-H_AD^2\mathcal F(A)[v,v]$. Moreover,
$H_A^TE(A)=m(A,B)$ by \eqref{nm-flat-identities} and
\eqref{nm-multiplier}. The exact reduced Hessian is therefore
\begin{equation}\label{nm-reduced-hessian}
 \begin{split}
 D_\xi^2[\mathcal S(A(\xi,B))][z,z]
 ={}&\langle v,D^2\mathcal S(A)v\rangle\\
 &-\langle m(A,B),D^2\mathcal F(A)[v,v]\rangle.
 \end{split}
\end{equation}
By \eqref{nm-action-hessian}, \eqref{nm-tangent-energy}, and
\eqref{nm-tangent-hilbert}, the first term is at least
$(1-Cr)\langle z,Kz\rangle$. Equations
\eqref{nm-hessian-composite}, \eqref{nm-multiplier-bounds}, and
\eqref{nm-tangent-hilbert} bound the second term in absolute value
by $C(b+r^2)\langle z,Kz\rangle$. Thus throughout a sufficiently
small real chart,
\begin{equation}\label{nm-strict-hessian}
 D_\xi^2[\mathcal S(A(\xi,B))][z,z]
 \ge\tfrac12\langle z,Kz\rangle
 \ge c\norm{z}_{2,\eta}^2.
\end{equation}

\begin{theorem}[Uniform local minimum on the fixed flat gauge slice]
\label{nm-theorem}
For sufficiently small real coarse $B$, the analytic field $A(B)$
constructed above is the unique stationary point in a fixed small
chart of
\[
 \mathcal F(A)=B,\qquad RD^TA=0.
\]
It is a strict local minimum of the normalized Wilson action on
that chart. Its radius, the bounds \eqref{nm-quadratic-remainder},
the analytic derivative estimates \eqref{nm-cauchy}, and the reduced
Hessian bound \eqref{nm-strict-hessian} are independent of
refinement and compatible periodic volume.
\end{theorem}

\begin{proof}
The contraction supplies the unique small stationary field. The
converse chart following \eqref{nm-chart} includes every small
feasible field in the fixed gauge slice. Integrating the positive
Hessian \eqref{nm-strict-hessian} along a line segment in the real
$\xi$-chart proves strict local minimality and local uniqueness.
\end{proof}

\subsection{Scope and normalization}
\label{nm-scope}

The statement concerns the corrected fixed-reference constraint
and a fixed flat gauge condition. A separate gauge-orbit argument
is required to transfer it to an unrestricted constrained minimum
or to a constraint formulated with ordinary block averages. It
does not assert a minimum among large fields or a nonflat source
gauge identification.

In dimension four the coupling-free action used here agrees with
the deposited Wilson normalization: Raw I:464--479 puts $g^{-2}$
outside the sum of $1-\operatorname{Re}\tau(P)$, and
I:12756--12760 specifies the physical trace pairing. The source
left-exponential Hessian at a nonzero background is a different
coordinate Hessian from the ordinary-potential Hessian used here;
the latter has been derived directly from the same action.

On the infinite lattice the same uniformly local maps can define
bounded stationary fields, but bounded fields generally have
infinite total action. Without further finite-energy hypotheses,
one must formulate stationarity and admissible local or decaying
comparisons rather than a minimum of that total action. This
module makes its minimum claim on finite periodic regulators.

\end{document}
